\subsection{Computation of the connecting homomorphisms}

Consider a commutative diagram

\begin{enumerate}
\def\labelenumi{(\arabic{enumi})}
\tightlist
\item
\end{enumerate}

\[
\begin{array}{c} 0 \longrightarrow \mathrm{R} ^ {\prime} \xrightarrow {\tilde {u}} \mathrm{R} \xrightarrow {\tilde {v}} \mathrm{R} ^ {\prime \prime} \longrightarrow 0 \\ a ^ {\prime} \Big \downarrow \qquad \qquad \qquad \qquad \qquad \qquad \qquad \qquad \qquad \qquad \qquad \qquad \qquad \qquad \qquad \qquad \qquad \qquad \qquad \qquad \qquad \qquad \qquad \qquad \qquad \qquad \qquad \qquad \qquad \qquad \qquad \qquad \qquad \qquad \\ 0 \longrightarrow \mathrm{M} ^ {\prime} \xrightarrow {u} \mathrm{M} \xrightarrow {v} \mathrm{M} ^ {\prime \prime} \longrightarrow 0 \end{array}\tag{2}
\]

where the first row (1) is an exact sequence of complexes of left A-modules, the second row (2) is an exact sequence of left A-modules, and where the vertical arrows are left resolutions.

PROPOSITION 3. --- a) Let P be an A-module, \(b: S \to P\) a left resolution of P; suppose that the sequence of complexes of k-modules

\[
0 \rightarrow \mathrm{S} \otimes_ {\mathrm{A}} \mathrm{R} ^ {\prime} \xrightarrow {l _ {\mathrm{s}} \otimes \tilde {u}} \mathrm{S} \otimes_ {\mathrm{A}} \mathrm{R} \xrightarrow {l _ {\mathrm{s}} \otimes \tilde {v}} \mathrm{S} \otimes_ {\mathrm{A}} \mathrm{R} ^ {\prime \prime} \rightarrow 0\tag{3}
\]

be exact. Then the following diagram is commutative:

\[
\begin{array}{c c c} \operatorname{Tor} ^ {\mathbf {A}} (\mathrm{P}, \mathrm{M} ^ {\prime \prime}) & \xrightarrow {\partial (\mathrm{P} , (2))} & \operatorname{Tor} ^ {\mathbf {A}} (\mathrm{P}, \mathrm{M} ^ {\prime}) \\ \psi (\mathrm{S}, \mathrm{R} ^ {\prime \prime}) \Big \downarrow & & \Big \downarrow \psi (\mathrm{S}, \mathrm{R} ^ {\prime}) \\ \mathrm{H} (\mathrm{S} \otimes_ {\mathbf {A}} \mathrm{R} ^ {\prime \prime}) & \xrightarrow {\partial ((3))} & \mathrm{H} (\mathrm{S} \otimes_ {\mathbf {A}} \mathrm{R} ^ {\prime}). \end{array}
\]

\begin{enumerate}
\def\labelenumi{\alph{enumi})}
\setcounter{enumi}{1}
\tightlist
\item
  Let N be a left A-module, \(c : N \to E\) a right resolution of N; suppose that the sequence of complexes of k-modules
\end{enumerate}

\begin{enumerate}
\def\labelenumi{(\arabic{enumi})}
\setcounter{enumi}{3}
\tightlist
\item
\end{enumerate}

\(0 \to \text{Homgr}_{\mathbf{A}}(\mathbf{R}^{\prime \prime}, \mathbf{E}) \xrightarrow{\text{Homgr}_{\mathbf{A}}(\tilde{\imath}, 1)} \text{Homgr}_{\mathbf{A}}(\mathbf{R}, \mathbf{E}) \xrightarrow{\text{Homgr}_{\mathbf{A}}(\tilde{\mu}, 1)} \text{Homgr}_{\mathbf{A}}(\mathbf{R}^{\prime}, \mathbf{E}) \to 0\) be exact. Then the following diagram is commutative:

\[
\begin{array}{c} \mathrm{H} (\operatorname{Homgr} _ {\mathbf {A}} (\mathrm{R} ^ {\prime}, \mathrm{E})) \xrightarrow {\delta ((4))} \mathrm{H} (\operatorname{Homgr} _ {\mathbf {A}} (\mathrm{R} ^ {\prime \prime}, \mathrm{E})) \\ \varphi (\mathrm{R} ^ {\prime}, \mathrm{E}) \Big \downarrow \\ \operatorname{Ext} _ {\mathbf {A}} (\mathrm{M} ^ {\prime}, \mathrm{N}) \xrightarrow {\partial ((2) , \mathrm{N})} \operatorname{Ext} _ {\mathbf {A}} (\mathrm{M} ^ {\prime \prime}, \mathrm{N}). \end{array}
\]

Let us prove for example \(a\) ). Let \(\beta: L(P) \to S\) be a morphism of complexes such that \(b \circ \beta = p_{P}\) ; consider the diagram of \(k\) -complexes

\[
\begin{array}{c}0 \rightarrow \mathrm{S} \otimes_ {\mathrm{A}} \mathrm{R} ^ {\prime} \xrightarrow {1 \otimes \tilde {u}} \mathrm{S} \otimes_ {\mathrm{A}} \mathrm{R} \xrightarrow {1 \otimes \tilde {v}} \mathrm{S} \otimes_ {\mathrm{A}} \mathrm{R} ^ {\prime \prime} \longrightarrow 0\\\beta \otimes 1 _ {\mathrm{R}} \Bigg | \quad \beta \otimes 1 _ {\mathrm{R}} \Bigg | \quad \beta \otimes 1 _ {\mathrm{R}} \Bigg |\\0 \rightarrow \mathrm{L(P)} \otimes_ {\mathrm{A}} \mathrm{R} ^ {\prime} \xrightarrow {1 \otimes \tilde {u}} \mathrm{L(P)} \otimes_ {\mathrm{A}} \mathrm{R} \xrightarrow {1 \otimes \tilde {v}} \mathrm{L(P)} \otimes_ {\mathrm{A}} \mathrm{R} ^ {\prime \prime} \longrightarrow 0\\1 \otimes a ^ {\prime} \Bigg | \quad 1 \otimes a \Bigg | \quad 1 \otimes a ^ {\prime \prime} \Bigg |\\0 \rightarrow \mathrm{L(P)} \otimes_ {\mathrm{A}} \mathrm{M} ^ {\prime} \xrightarrow {1 \otimes u} \mathrm{L(P)} \otimes_ {\mathrm{A}} \mathrm{M} \xrightarrow {1 \otimes v} \mathrm{L(P)} \otimes_ {\mathrm{A}} \mathrm{M} ^ {\prime \prime} \longrightarrow 0.\end{array}
\]

It is commutative with exact rows (by the hypothesis for the first row, and by the fact that L(P) is flat for the other two). We therefore have a commutative diagram (X, p.~31, prop. 2, and X, p.~72, def. 2).

\[
\begin{array}{c c c} \mathrm{H} (\mathrm{S} \otimes_ {\mathrm{A}} \mathrm{R} ^ {\prime \prime}) & \xrightarrow {\partial ((3))} & \mathrm{H} (\mathrm{S} \otimes_ {\mathrm{A}} \mathrm{R} ^ {\prime}) \\ \mathrm{H} (\beta \otimes 1) \uparrow & & \uparrow \mathrm{H} (\beta \otimes 1) \\ \mathrm{H} (\mathrm{L} (\mathrm{P}) \otimes_ {\mathrm{A}} \mathrm{R} ^ {\prime \prime}) & \longrightarrow & \mathrm{H} (\mathrm{L} (\mathrm{P}) \otimes_ {\mathrm{A}} \mathrm{R} ^ {\prime}) \\ \mathrm{H} (1 \otimes a ^ {\prime \prime}) \downarrow & & \downarrow \mathrm{H} (1 \otimes a ^ {\prime}) \\ \mathrm{H} (\mathrm{L} (\mathrm{P}) \otimes_ {\mathrm{A}} \mathrm{M} ^ {\prime \prime}) & \longrightarrow & \mathrm{H} (\mathrm{L} (\mathrm{P}) \otimes_ {\mathrm{A}} \mathrm{M} ^ {\prime}) \\ \psi_ {\mathrm{P}} (\mathbf {M} ^ {\prime \prime}) \uparrow & & \uparrow \psi_ {\mathrm{P}} (\mathbf {M} ^ {\prime}) \\ \operatorname{Tor} (\mathrm{P}, \mathbf {M} ^ {\prime \prime}) & \xrightarrow {\partial (\mathrm{P} , (2))} & \operatorname{Tor} (\mathrm{P}, \mathbf {M} ^ {\prime}). \end{array}
\]

By X, p.~67, prop. 4, H(1 ⊗ a′′) and H(1 ⊗ a′) are bijective; on the other hand, by definition of the homomorphisms ψ, we have H(β ⊗ 1) ∘ ψ(L(P), R′′) = ψ(S, R′′) and H(1 ⊗ a′′)∘ψ(L(P), R′′) = ψ(L(P), M′′) = ψ\_P(M′′), hence

\[
\psi (\mathrm{S}, \mathrm{R} ^ {\prime \prime}) = \mathrm{H} (\beta \otimes 1) \circ \mathrm{H} (1 \otimes a ^ {\prime \prime}) ^ {- 1} \circ \psi_ {\mathrm{P}} (\mathrm{M} ^ {\prime \prime});
\]

likewise, \(\psi(S, R') = H(\beta \otimes 1) \circ H(1 \otimes a')^{-1} \circ \psi_P(M')\) , and the desired assertion \(\partial((3)) \circ \psi(S, R'') = \psi(S, R') \circ \partial(P, (2))\) follows from the commutativity of the preceding diagram.

Remarks. --- 1) Using the commutation isomorphisms, one deduces from a) the analogous statement obtained by interchanging the roles of the two arguments of the tensor product.

\begin{enumerate}
\def\labelenumi{\arabic{enumi})}
\setcounter{enumi}{1}
\tightlist
\item
  With the notations of \(a\) ), suppose either S is flat, or R, R', R'' flat; then on the one hand the sequence (3) is exact (X, p.~72, cor. 2) and one can apply Prop. 3; on the other hand \(\psi(S, R')\) is bijective (th. 1), hence
\end{enumerate}

\[
\partial (\mathrm{P}, (2)) = \psi (\mathrm{S}, \mathrm{R} ^ {\prime}) ^ {- 1} \circ \partial ((3)) \circ \psi (\mathrm{S}, \mathrm{R} ^ {\prime \prime}).
\]

\begin{enumerate}
\def\labelenumi{\arabic{enumi})}
\setcounter{enumi}{2}
\tightlist
\item
  With the notations of \(b\) ), suppose either E is injective, or R, R', R'\,' projective; then on the one hand the sequence (4) is exact (X, p.~83, prop. 2) and one can apply prop. 3; on the other hand, \(\varphi(R', E)\) is bijective (th. 1); therefore
\end{enumerate}

\[
\delta ((2), \mathrm{N}) = \varphi (\mathrm{R} ^ {\prime \prime}, \mathrm{E}) \circ \hat {\partial} ((4)) \circ \varphi (\mathrm{R} ^ {\prime}, \mathrm{E}) ^ {- 1}.
\]

Consider now a commutative diagram

\begin{enumerate}
\def\labelenumi{(\arabic{enumi})}
\setcounter{enumi}{4}
\tightlist
\item
\end{enumerate}

\[
\begin{array}{c} 0 \longrightarrow \mathrm{N} ^ {\prime} \xrightarrow {r} \mathrm{N} \xrightarrow {s} \mathrm{N} ^ {\prime \prime} \longrightarrow 0 \\ c ^ {\prime} \Big \downarrow \qquad \qquad c \Big \downarrow \qquad \qquad c ^ {\prime \prime} \Big \downarrow \\ 0 \longrightarrow \mathrm{E} ^ {\prime} \xrightarrow {\tilde {r}} \mathrm{E} \xrightarrow {\tilde {s}} \mathrm{E} ^ {\prime \prime} \longrightarrow 0 \end{array}\tag{6}
\]

whose first row (5) is an exact sequence of left A-modules, the second row (6) an exact sequence of complexes of left A-modules, and where the vertical arrows are right resolutions. Reasoning as in Prop. 3, one proves the following proposition:

PROPOSITION 4. --- Let M be a left A-module, \(a: \mathbf{R} \to \mathbf{M}\) a left resolution of M such that the sequence

\begin{enumerate}
\def\labelenumi{(\arabic{enumi})}
\setcounter{enumi}{6}
\tightlist
\item
  \(0 \to \operatorname{Homgr}_{\mathbf{A}}(\mathbf{R}, \mathbf{E}') \xrightarrow{\operatorname{Homgr}(\tilde{\mathcal{S}}, 1)} \operatorname{Homgr}_{\mathbf{A}}(\mathbf{R}, \mathbf{E}) \xrightarrow{\operatorname{Homgr}(\tilde{\mathcal{F}}, 1)} \operatorname{Homgr}_{\mathbf{A}}(\mathbf{R}, \mathbf{E}'' \to 0\) is exact. Then the following diagram is commutative.
\end{enumerate}

\[
\begin{array}{c} \mathrm{H} (\operatorname{Homgr} _ {\mathbf {A}} (\mathbf {R}, \mathbf {E} ^ {\prime \prime})) \xrightarrow {\partial ((7))} \mathrm{H} (\operatorname{Homgr} _ {\mathbf {A}} (\mathbf {R}, \mathbf {E} ^ {\prime})) \\ \varphi (\mathbf {R}, \mathbf {E} ^ {\prime \prime}) \Biggl \downarrow \\ \operatorname{Ext} _ {\mathbf {A}} (\mathbf {M}, \mathbf {N} ^ {\prime \prime}) \xrightarrow {\delta (\mathbf {M} , (5))} \operatorname{Ext} _ {\mathbf {A}} (\mathbf {M}, \mathbf {N} ^ {\prime}). \end{array}
\]

Remarks. --- 4) If R is projective or if E, E', E'' are injective, the sequence (7) is exact (X, p.~83, prop. 2); moreover, φ(R, E'') is bijective (th. 1); hence

\[
\delta (M, (5)) = \varphi (R, E ^ {\prime}) \circ \partial ((7)) \circ \varphi (R, E ^ {\prime \prime}) ^ {- 1}.
\]

\begin{enumerate}
\def\labelenumi{\arabic{enumi})}
\setcounter{enumi}{4}
\tightlist
\item
  We leave it to the reader to state and prove the propositions analogous to Props. 3 and 4 and relating to homomorphisms. \(\psi_{1}\) and \(\varphi_{1}\) .
\end{enumerate}

